% Options for packages loaded elsewhere
\PassOptionsToPackage{unicode}{hyperref}
\PassOptionsToPackage{hyphens}{url}
\PassOptionsToPackage{dvipsnames,svgnames,x11names}{xcolor}
%
\documentclass[
  10pt,
  a4paper,
]{article}
\usepackage{amsmath,amssymb}
\usepackage{iftex}
\ifPDFTeX
  \usepackage[T1]{fontenc}
  \usepackage[utf8]{inputenc}
  \usepackage{textcomp} % provide euro and other symbols
\else % if luatex or xetex
  \usepackage{unicode-math} % this also loads fontspec
  \defaultfontfeatures{Scale=MatchLowercase}
  \defaultfontfeatures[\rmfamily]{Ligatures=TeX,Scale=1}
\fi
\usepackage{lmodern}
\ifPDFTeX\else
  % xetex/luatex font selection
\fi
% Use upquote if available, for straight quotes in verbatim environments
\IfFileExists{upquote.sty}{\usepackage{upquote}}{}
\IfFileExists{microtype.sty}{% use microtype if available
  \usepackage[]{microtype}
  \UseMicrotypeSet[protrusion]{basicmath} % disable protrusion for tt fonts
}{}
\makeatletter
\@ifundefined{KOMAClassName}{% if non-KOMA class
  \IfFileExists{parskip.sty}{%
    \usepackage{parskip}
  }{% else
    \setlength{\parindent}{0pt}
    \setlength{\parskip}{6pt plus 2pt minus 1pt}}
}{% if KOMA class
  \KOMAoptions{parskip=half}}
\makeatother
\usepackage{xcolor}
\usepackage[margin=1in]{geometry}
\usepackage{color}
\usepackage{fancyvrb}
\newcommand{\VerbBar}{|}
\newcommand{\VERB}{\Verb[commandchars=\\\{\}]}
\DefineVerbatimEnvironment{Highlighting}{Verbatim}{commandchars=\\\{\}}
% Add ',fontsize=\small' for more characters per line
\newenvironment{Shaded}{}{}
\newcommand{\AlertTok}[1]{\textcolor[rgb]{1.00,0.00,0.00}{\textbf{#1}}}
\newcommand{\AnnotationTok}[1]{\textcolor[rgb]{0.38,0.63,0.69}{\textbf{\textit{#1}}}}
\newcommand{\AttributeTok}[1]{\textcolor[rgb]{0.49,0.56,0.16}{#1}}
\newcommand{\BaseNTok}[1]{\textcolor[rgb]{0.25,0.63,0.44}{#1}}
\newcommand{\BuiltInTok}[1]{\textcolor[rgb]{0.00,0.50,0.00}{#1}}
\newcommand{\CharTok}[1]{\textcolor[rgb]{0.25,0.44,0.63}{#1}}
\newcommand{\CommentTok}[1]{\textcolor[rgb]{0.38,0.63,0.69}{\textit{#1}}}
\newcommand{\CommentVarTok}[1]{\textcolor[rgb]{0.38,0.63,0.69}{\textbf{\textit{#1}}}}
\newcommand{\ConstantTok}[1]{\textcolor[rgb]{0.53,0.00,0.00}{#1}}
\newcommand{\ControlFlowTok}[1]{\textcolor[rgb]{0.00,0.44,0.13}{\textbf{#1}}}
\newcommand{\DataTypeTok}[1]{\textcolor[rgb]{0.56,0.13,0.00}{#1}}
\newcommand{\DecValTok}[1]{\textcolor[rgb]{0.25,0.63,0.44}{#1}}
\newcommand{\DocumentationTok}[1]{\textcolor[rgb]{0.73,0.13,0.13}{\textit{#1}}}
\newcommand{\ErrorTok}[1]{\textcolor[rgb]{1.00,0.00,0.00}{\textbf{#1}}}
\newcommand{\ExtensionTok}[1]{#1}
\newcommand{\FloatTok}[1]{\textcolor[rgb]{0.25,0.63,0.44}{#1}}
\newcommand{\FunctionTok}[1]{\textcolor[rgb]{0.02,0.16,0.49}{#1}}
\newcommand{\ImportTok}[1]{\textcolor[rgb]{0.00,0.50,0.00}{\textbf{#1}}}
\newcommand{\InformationTok}[1]{\textcolor[rgb]{0.38,0.63,0.69}{\textbf{\textit{#1}}}}
\newcommand{\KeywordTok}[1]{\textcolor[rgb]{0.00,0.44,0.13}{\textbf{#1}}}
\newcommand{\NormalTok}[1]{#1}
\newcommand{\OperatorTok}[1]{\textcolor[rgb]{0.40,0.40,0.40}{#1}}
\newcommand{\OtherTok}[1]{\textcolor[rgb]{0.00,0.44,0.13}{#1}}
\newcommand{\PreprocessorTok}[1]{\textcolor[rgb]{0.74,0.48,0.00}{#1}}
\newcommand{\RegionMarkerTok}[1]{#1}
\newcommand{\SpecialCharTok}[1]{\textcolor[rgb]{0.25,0.44,0.63}{#1}}
\newcommand{\SpecialStringTok}[1]{\textcolor[rgb]{0.73,0.40,0.53}{#1}}
\newcommand{\StringTok}[1]{\textcolor[rgb]{0.25,0.44,0.63}{#1}}
\newcommand{\VariableTok}[1]{\textcolor[rgb]{0.10,0.09,0.49}{#1}}
\newcommand{\VerbatimStringTok}[1]{\textcolor[rgb]{0.25,0.44,0.63}{#1}}
\newcommand{\WarningTok}[1]{\textcolor[rgb]{0.38,0.63,0.69}{\textbf{\textit{#1}}}}
\setlength{\emergencystretch}{3em} % prevent overfull lines
\providecommand{\tightlist}{%
  \setlength{\itemsep}{0pt}\setlength{\parskip}{0pt}}
\setcounter{secnumdepth}{-\maxdimen} % remove section numbering
\setlength{\emergencystretch}{3em}
\ifLuaTeX
  \usepackage{selnolig}  % disable illegal ligatures
\fi
\usepackage{bookmark}
\IfFileExists{xurl.sty}{\usepackage{xurl}}{} % add URL line breaks if available
\urlstyle{same}
\hypersetup{
  colorlinks=true,
  linkcolor={blue},
  filecolor={Maroon},
  citecolor={Blue},
  urlcolor={blue},
  pdfcreator={LaTeX via pandoc}}

\author{}
\date{}

\begin{document}

\section{DDTA minimal BAE identity ontology -
R1}\label{ddta-minimal-bae-identity-ontology---r1}

\textbf{Status:} CLOSED BY BA1-T3
\textbf{R24 alignment note:} BA1 identity semantics remain retained. Statements below that describe BA2 as \texttt{NOT STARTED / NEXT AUTHORIZED PHASE} are closure-time sequencing, not current execution state. Current R24 forward state is recorded in \nolinkurl{DDTA_R24_DECISION_RULE_CHECKPOINT.md}; the active BA2 working revision is \nolinkurl{BA2_RELATION_ACTION_VOCABULARY_R2.md}.


\textbf{Repository baseline reviewed:}
\texttt{e88d7e220536863d564f9e3b9fac7f1592a8c440}

\textbf{Activation:} applies when the BA1-T3 closure package is accepted
and committed.

\textbf{Authority note:} this artifact defines the minimal Base Analysis
semantic identity boundary. Governed project documentation remains
project authority.

\subsection{1. Closed ontology
statement}\label{closed-ontology-statement}

For the current thesis scope, the minimal first-class semantic identity
ontology of Base Analysis contains exactly two families:

\begin{Shaded}
\begin{Highlighting}[]
\NormalTok{BAReferent}
\NormalTok{BAProposition}
\end{Highlighting}
\end{Shaded}

\texttt{BAE} is an umbrella term for Base Analysis elements. BA1 does
not require an additional \texttt{BAE} metaclass.

The two families are semantic distinctions. One implementation may
materialize them in separate classes/tables, in one physical container
with a discriminator, or in another representation, provided the closed
semantic distinction is preserved.

\subsection{\texorpdfstring{2. \texttt{BAReferent} -
ACCEPTED}{2. BAReferent - ACCEPTED}}\label{bareferent---accepted}

A \texttt{BAReferent} is an independently identifiable unit of
methodology-neutral shared project meaning that must be reusable across
one or more analytical propositions, projections or governed baselines.

A BAReferent is not restricted to a concrete runtime object. Depending
on governed project meaning, it may denote a participant, capability,
information/resource concept, behavior/event, milestone, contract,
store, boundary, state/context or another independently reusable
semantic referent.

When project meaning must itself be reused, constrained, correlated,
compared across baselines or targeted by multiple assertions, that
meaning receives BAReferent identity.

\subsection{\texorpdfstring{3. \texttt{BAProposition} -
ACCEPTED}{3. BAProposition - ACCEPTED}}\label{baproposition---accepted}

A \texttt{BAProposition} is an independently identifiable
methodology-neutral analytical assertion about one or more BAReferent
identities, with enough independent identity to support assertion-level
origin/provenance, diagnosis, reuse and change disposition.

A BAProposition is not the project-semantic thing it describes. If the
underlying project meaning needs independent reuse or qualification,
that meaning is represented as a BAReferent and propositions state the
relevant semantics about it.

The exact internal structure of a proposition, its participation roles,
arity, predicates, actions and cardinalities are not closed by BA1. They
belong to BA2.

\subsection{4. Why two identity families are
necessary}\label{why-two-identity-families-are-necessary}

Removing BAReferent would force propositions to point to raw source
strings, document identities or ad-hoc duplicated concepts, weakening
stable shared semantic identity, projection reuse and change-impact
reasoning.

Removing BAProposition would force assertion-level origin, uncertainty
and change disposition into properties of referents or into anonymous
relation rows. The reviewed corpora require the same referent to survive
while assertions about it are introduced, revised, superseded, retained
or diagnosed. Reifying every assertion as an undifferentiated referent
would merely rename proposition semantics rather than eliminate them.

Therefore semantic collapse into one undifferentiated family is
rejected.

\subsection{5. Domain kinds are not automatically first-class
types}\label{domain-kinds-are-not-automatically-first-class-types}

BA1 does not require dedicated first-class metaclasses for:

\begin{itemize}
\tightlist
\item
  Participant / Actor;
\item
  Component / Capability;
\item
  Information / Resource;
\item
  Behavior / Event / Transition;
\item
  Interface / Connection / Channel;
\item
  Store / Persistence;
\item
  Contract;
\item
  Boundary / Externality;
\item
  State / Mode / Context;
\item
  Dependency / Allocation;
\item
  Property / Constraint;
\item
  DataFlow / Interaction.
\end{itemize}

A concrete meaning from any of these categories may still be represented
as a BAReferent when independent identity is required.

Its category may be preserved through method-neutral classification,
roles, values or propositions. BA1 intentionally does not freeze the
exact classification vocabulary or cardinality.

\subsection{6. First-class split criterion -
CLOSED}\label{first-class-split-criterion---closed}

A future dedicated BAE type/family is justified only when concrete
evidence shows both:

\begin{enumerate}
\def\labelenumi{\arabic{enumi}.}
\tightlist
\item
  \textbf{independent semantic identity} is required across
  propositions, projections or change; and
\item
  \textbf{reusable subtype-specific invariants} cannot be represented
  honestly as classifications, roles, values or propositions over the
  accepted families.
\end{enumerate}

Recurring usefulness, familiar systems-modeling terminology, diagram
convenience, programming classes, database tables or a method-specific
category are insufficient by themselves.

\subsection{7. Closed family-boundary
invariants}\label{closed-family-boundary-invariants}

\begin{enumerate}
\def\labelenumi{\arabic{enumi}.}
\tightlist
\item
  \textbf{Two-family minimum:} Base Analysis preserves the semantic
  distinction between project-semantic referent identity and analytical
  assertion identity.
\item
  \textbf{Referent reuse:} independently reusable project meaning is
  represented through BAReferent identity.
\item
  \textbf{Assertion identity:} independently
  reviewable/provenance-bearing shared analytical claims are represented
  through BAProposition identity.
\item
  \textbf{No proposition proxy:} BAProposition is not required to serve
  as a proxy project referent; reusable project meaning receives
  BAReferent identity.
\item
  \textbf{Method neutrality:} neither family embeds STRIDE, STRIDE-AI,
  DFD or tool-owned semantic categories.
\item
  \textbf{Document separation:} governed document kinds remain
  source-layer concepts and are not copied into Base Analysis as BAE
  types by inheritance.
\item
  \textbf{Classification is cheaper than splitting:} method-neutral
  semantic kind or participation-role information does not create a new
  first-class type unless the split criterion is met.
\item
  \textbf{Representation independence:} the semantic family boundary
  does not mandate a programming class hierarchy, graph notation,
  database schema or persisted canonical graph.
\item
  \textbf{Authority preservation:} BAReferent and BAProposition
  identities do not make Base Analysis project authority; governed
  documentation remains authoritative.
\item
  \textbf{Phase separation:} relation/action structure is BA2;
  provenance/identity mechanics are BA3; projections BA4; vocabulary
  assistance BA5; implemented Base Analysis regression BA6.
\end{enumerate}

\subsection{8. Explicit non-types and downstream
constructs}\label{explicit-non-types-and-downstream-constructs}

The following are not BA1 domain BAE metaclasses:

\begin{itemize}
\tightlist
\item
  document-layer roots/commitments: \texttt{MacroRequirement} and
  \texttt{Decision};
\item
  document-layer requirement types: \texttt{Requirement} and
  \texttt{FunctionalRequirement};
\item
  \texttt{SpecializedRequirement} (document layer);
\item
  \texttt{SecurityRequirement} and \texttt{NormativeClause} (document
  layer);
\item
  grounded, derived and diagnostic/unresolved origin/state categories;
\item
  source locator or provenance record;
\item
  view/projection;
\item
  AnalysisRecord, Finding or Common Finding;
\item
  ThreatForge implementation objects;
\item
  STRIDE/STRIDE-AI/DFD elements.
\end{itemize}

Their relevant responsibilities remain source-layer, metadata,
projection, analysis-envelope or later-phase concerns as already
allocated by BA0 and the research work plan.

\subsection{9. Correct interpretation of
specialization}\label{correct-interpretation-of-specialization}

BA1 closure does not claim that all BAReferents are semantically
identical.

A facial-access delivery behavior, a RecognitionCapture information
concept, a PaymentProviderContract and an InventoryLedger can all be
distinct BAReferents with different method-neutral classifications and
different propositions.

The closure only says that current evidence does not require separate
\textbf{first-class identity families} for those categories. BA2 may
define roles/relations/classification mechanics needed to distinguish
and consume them. A later concrete counterexample may reopen BA1 if
classification is insufficient and a true subtype-specific invariant is
required.

\subsection{10. Reopen criteria}\label{reopen-criteria}

Reopen BA1 only for a material identity-ontology counterexample, for
example:

\begin{enumerate}
\def\labelenumi{\arabic{enumi}.}
\tightlist
\item
  a governed corpus requires a semantic kind with independent identity
  plus subtype-specific invariants that cannot be represented honestly
  over BAReferent/BAProposition;
\item
  a method-neutral consumer cannot construct its required projection
  without independently reconstructing a missing first-class semantic
  distinction from raw prose;
\item
  BA2 cannot define coherent proposition participation/role semantics
  without introducing an unacknowledged third identity family;
\item
  BA3 shows that assertion-level provenance/diagnosis/change cannot be
  attached coherently to BAProposition identity;
\item
  a concrete case proves BAProposition identity can be removed without
  losing any BA0 responsibility, materially falsifying the two-family
  lower bound.
\end{enumerate}

Do not reopen BA1 merely because a dedicated subtype would make code, a
diagram, a query or one analysis method more convenient.

\subsection{11. Closure disposition}\label{closure-disposition}

\begin{Shaded}
\begin{Highlighting}[]
\NormalTok{BA1             CLOSED}
\NormalTok{BAReferent      ACCEPTED first{-}class semantic identity family}
\NormalTok{BAProposition   ACCEPTED first{-}class semantic identity family}
\NormalTok{BAE             umbrella term; no additional metaclass required}
\NormalTok{Domain splits   NOT REQUIRED by current evidence}
\NormalTok{BA2             NOT STARTED / NEXT AUTHORIZED PHASE}
\end{Highlighting}
\end{Shaded}


\end{document}
